\subsection{疲劳分析}
\label{sec:fatigue-filter}

\subsubsection{取值}

多功能过滤工具的五个承力件（提升短节、滤进接头、加长过滤套筒、滤出接头、下接头）均为 4145H 调质钢，图纸技术要求 $\sigma_s\ge930$\,MPa、$\sigma_u\ge1\,080$\,MPa。该档材料与 4330V 不同，S-N 曲线须单独建立：$k_a=4.51\times1\,080^{-0.265}=0.709$（$Ra1.6$ 机加工面）、$k_b=0.748$、$k_c=0.85$、$k_d=1.0$、$k_e=0.814$，五项系数之积为 0.367，得 $1\times10^{6}$ 次处 $S_e=0.5\times1\,080\times0.367\approx198.1$\,MPa。

$1\times10^{3}$ 次处的取值需额外处理一步：$0.9\sigma_u=972$\,MPa 已超过屈服强度 930\,MPa，若照 $0.9\sigma_u$ 取锚点，曲线在低周段将给出材料无法承受的应力幅。按第~\ref{ss:fatigue-probe}~节的规则以屈服强度截断，取 930\,MPa，五点拟合得 $\sigma_a=4\,366\,N^{-0.2239}$，$1\times10^{3}$ 次处许用应力幅 $\sigma_e=930$\,MPa。截断使曲线在低周段较幂律延伸略平，本工具的评定点落于 $10^{4}$ 次量级，处于拟合段内，不受截断方式的影响。

轴向拉伸工况的峰值 980.45\,MPa 位于加长过滤套筒与滤出接头连接处的变径台阶根部，属几何尖角引起的数值奇异。同一位置附近的三个探针读数为 600.56\,MPa（节点 176608）、632.05\,MPa 与 742.11\,MPa，彼此相差 24\%，表明该处应力梯度很陡；取梯度已平缓的 600.56\,MPa 作为评定应力，与前期静力复核的 597\,MPa 相差 0.6\%，两种模型相互印证。抗扭工况按 18\,kN$\cdot$m 施加（该工具合同扭矩为 18\,kN$\cdot$m，高于其余工具的 17\,kN$\cdot$m），峰值 832.37\,MPa 位于同一位置，评定应力为 527.75\,MPa（节点 61155），较轴向工况低 12\%，与静力复核的 544\,MPa 相差 3\%。扭矩按完全反向 $R=-1$ 取定，其理由与强磁工具相同：扭矩会换向，且正反转与憋钻释放时冲击最大。评定结果列于表~\ref{tab:fatigue_filter}。

\begin{table}[htbp]
\centering
\caption{多功能过滤工具的疲劳评定结果（200 趟）}
\label{tab:fatigue_filter}
\begin{tabularx}{\textwidth}{lCCC}
\hline
项目 & 轴向拉伸（$R=0$） & 抗扭（$R=-1$） & 判据 \\
\hline
额定工况 & 1\,000\,kN & 18\,kN$\cdot$m & --- \\
评定应力 / MPa & 600.6 & 527.8 & --- \\
应力幅与平均应力 / MPa & 300.3 & 527.8 / 0 & --- \\
等效对称应力幅 / MPa & 413.9 & 527.8 & --- \\
计算寿命 / 次 & $3.7\times10^{4}$ & $1.25\times10^{4}$ & --- \\
200 次累积损伤 & $5.4\times10^{-3}$ & $1.6\times10^{-2}$ & $\le 1.0$ \\
疲劳安全系数 & 1.66 & 1.76 & $\ge 1.5$ \\
静强度安全系数 & 1.55 & 1.76 & $\ge 1.5$ \\
寿命裕度 & 约 185 & 约 63 & --- \\
\hline
\end{tabularx}
\end{table}

拉伸工况的计算流程与其余工具一致：$\sigma_a=\sigma_m=300.3$\,MPa，$\sigma_{a,\mathrm{eq}}=300.3/(1-300.3/1\,080)=413.9$\,MPa，按 4145H 的幂律反解得 $N=(4\,366/413.9)^{1/0.2239}=3.7\times10^{4}$ 次，$D=200/3.7\times10^{4}=5.4\times10^{-3}$，$n_f=1/(300.3/930+300.3/1\,080)=1.66$，$n_s=930/600.6=1.55$。抗扭工况 $\sigma_m=0$，$N=1.25\times10^{4}$ 次（云图探针读数 12\,526 次，二者一致），$D=1.6\times10^{-2}$，$n_f=930/527.8=1.76$，与 $n_s$ 相同。

两个工况均满足判据，但轴向工况的静强度安全系数 1.55 距判据 1.5 仅 3\%，与强磁工具同属裕度较紧的一类。尤应注意采样位置的影响：若以 742.11\,MPa 评定，静强度安全系数仅为 1.25，将不满足判据。本文按与其他工具一致的规则取 600.56\,MPa，并将 742.11\,MPa 单列备查，建议对该过渡圆角加密一档网格复核，同时将圆角半径加大至 R3 以上。加大圆角既可降低峰值，也可减小探针之间的梯度，使评定值对采样位置不再敏感。

抗扭工况较轴向工况宽裕：评定应力低 12\%，且因完全反向无需 Goodman 折算，寿命仍达 $1.25\times10^{4}$ 次；累积损伤 $1.6\times10^{-2}$ 较轴向工况的 $5.4\times10^{-3}$ 高一个量级，但远低于判据。

控制部位为上下接头与滤芯座的螺纹过渡段。除加工上应保证过渡圆角与螺纹根部光洁度、成品按 NB/T\,47013.4 进行磁粉探伤外，使用中还应注意：更换滤芯时不得用大锤敲击工具本体。敲击在螺纹过渡段引入的冲击应力远高于额定拉伸应力，是此类工具疲劳失效的常见诱因。
